\makeatletter
\renewcommand*\env@matrix[1][*\c@MaxMatrixCols c]{%
  \hskip -\arraycolsep
  \let\@ifnextchar\new@ifnextchar
  \array{#1}}
\makeatother

The constraint polyhedron is represented in the figure below:

\begin{center}
\begin{tikzpicture}[scale=1.4]
  \begin{axis}[ xlabel={$x_1$}, ylabel={$x_2$}
      ,axis on top=true
      ,axis equal
      ,axis lines=middle
      ,samples=41
      ,xtick = {1}
      ,ytick = {1}
      ,domain=0:1.5
      ,xmin=-1.3,xmax=1.3
      ,ymin=-0.1,ymax=1.2
,legend pos=south east
]
\addplot [thick,color=gray!20,fill=gray!20, 
                    fill opacity=0.05]coordinates {
  ( -1.3,0) 
  ( -1.3,1)
  ( 1.3,1)
  (1.3,0)
};
\addplot[color=black, thin,domain=-1.2:1.2] {0}
node[pos=0.2,sloped,above] {\tiny $x_2=0$};
\addplot[color=black,thin,domain=-1.2:1.2] {1}
node[pos=0.2,sloped,below] {\tiny$x_2=1$};

\node [circle,fill=black,draw,scale=0.2,label=above right:$A$] at (0,1) {A};

\end{axis}
\end{tikzpicture}
\end{center}

In order to write
the problem in standard form, we need to replace $x_1$ by the
difference of two non negative variables
\[
x_1 = x_1^+ - x_1^-,
\]
and we need to introduce a slack variable for the constraint $x_2 \leq 1$.
We obtain
\[
\min_{x \in \R^4} x_1^+ - x_1^-  + x_2
\]
subject to
\begin{align*}
  x_2 + x_3 &= 1, \\
  x_1^+, x_1^-, x_2, x_3 &\geq 0.
\end{align*}
In matrix form, we have
\[
\min_{x \in \R^4} c^T x
\]
subject to
\begin{align*}
  Ax &= b, \\
  x &\geq 0,
\end{align*}
where
\begin{align*}
  A &= \begin{pmatrix}
    0 & 0 & 1 & 1 
  \end{pmatrix}, \\
  b &= 1, \\
  c &= \rvect{1 \\ -1 \\ 1 \\ 0},\\
  x &= \cvect{x_1^+ \\ x_1^- \\ x_2 \\ x_3}.
\end{align*}
This is a problem with $n=4$ variables and $m=1$
constraints. So the basis matrix is of dimension one, and is a scalar.
If $x_2$ is in the basis, $B=1$ and we have
  \[
  x_2=B^{-1}b=1,
  \]
and the feasible basic solution is
  \[
\cvect{0 \\ 0 \\ 1 \\ 0},
\]
and corresponds to point A on the figure. As there are three non basic
variables, there are three basic directions.

\begin{enumerate}
\item The basic direction associated with $x_1^+$ is
  \[
d_1^+ = \cvect{1 \\ 0 \\ d_B \\ 0},
\]
where
\[
d_B = -B^{-1}A_1 = 0,
\]
as the variable $x_1^+$ is associated with column 1 of $A$,
so that 
  \[
d_1^+ = \cvect{1 \\ 0 \\ 0 \\ 0}.
\]
Performing a step $\alpha$ in this direction leads to the point
\[
\cvect{0 \\ 0 \\ 1 \\ 0} + \alpha \cvect{1 \\ 0 \\ 0 \\ 0} =
\cvect{\alpha \\ 0 \\ 1 \\ 0}.
\]
The point is feasible for any value of $\alpha \geq 0$, so that the
direction is feasible.
In the original variables, with $x_1=x_1^+-x_1^-$, it corresponds to
a direction
\[
\cvect{1-0 \\ 0} =
\cvect{1 \\ 0},
\]
as illustrated on the following figure:
\begin{center}
\begin{tikzpicture}[scale=1.2]
  \begin{axis}[ xlabel={$x_1$}, ylabel={$x_2$}
      ,axis on top=true
      ,axis equal
      ,axis lines=middle
      ,samples=41
      ,xtick = {1}
      ,ytick = {1}
      ,domain=0:1.5
      ,xmin=-1.3,xmax=1.3
      ,ymin=-0.1,ymax=1.2
,legend pos=south east
]
\addplot [thick,color=gray!20,fill=gray!20, 
                    fill opacity=0.05]coordinates {
  ( -1.3,0) 
  ( -1.3,1)
  ( 1.3,1)
  (1.3,0)
};
\addplot[color=black, thin,domain=-1.2:1.2] {0}
node[pos=0.2,sloped,above] {\tiny $x_2=0$};
\addplot[color=black,thin,domain=-1.2:1.2] {1}
node[pos=0.2,sloped,below] {\tiny$x_2=1$};

\draw[-latex,very thick] (0,1) -- (1,1) node[pos=0.5,above,sloped] {$d_1^+$} ;
%\node [circle,fill=black,draw,scale=0.2,label=above right:$A$] at (0,1) {A};
%\node [circle,fill=black,draw,scale=0.2,label=below right:$B$] at (0,0) {B};

\end{axis}
\end{tikzpicture}
\end{center}

\item The basic direction associated with $x_1^-$ is
  \[
d_1^- = \cvect{0 \\ 1 \\ d_B \\ 0},
\]
where
\[
d_B = -B^{-1}A_2 = 0,
\]
as the variable $x_1^-$ is associated with column 2 of $A$,
so that 
  \[
d_1^- = \cvect{0 \\ 1 \\ 0 \\ 0}.
\]
Performing a step $\alpha$ in this direction leads to the point
\[
\cvect{0 \\ 0 \\ 1 \\ 0} + \alpha \cvect{0 \\ 1 \\ 0 \\ 0} =
\cvect{0 \\ \alpha \\ 1 \\ 0}.
\]
The point is feasible for any value of $\alpha \geq 0$, so that the
direction is feasible.
In the original variables, with $x_1=x_1^+-x_1^-$, it corresponds to
a direction
\[
\cvect{0-1 \\ 0}=
\cvect{-1 \\ 0},
\]
as illustrated on the following figure:
\begin{center}
\begin{tikzpicture}[scale=1.2]
  \begin{axis}[ xlabel={$x_1$}, ylabel={$x_2$}
      ,axis on top=true
      ,axis equal
      ,axis lines=middle
      ,samples=41
      ,xtick = {1}
      ,ytick = {1}
      ,domain=0:1.5
      ,xmin=-1.3,xmax=1.3
      ,ymin=-0.1,ymax=1.2
,legend pos=south east
]
\addplot [thick,color=gray!20,fill=gray!20, 
                    fill opacity=0.05]coordinates {
  ( -1.3,0) 
  ( -1.3,1)
  ( 1.3,1)
  (1.3,0)
};
\addplot[color=black, thin,domain=-1.2:1.2] {0}
node[pos=0.2,sloped,above] {\tiny $x_2=0$};
\addplot[color=black,thin,domain=-1.2:1.2] {1}
node[pos=0.2,sloped,below] {\tiny$x_2=1$};

\draw[-latex,very thick] (0,1) -- (-1,1) node[pos=0.5,above,sloped] {$d_1^-$} ;
%\node [circle,fill=black,draw,scale=0.2,label=above right:$A$] at (0,1) {A};
%\node [circle,fill=black,draw,scale=0.2,label=below right:$B$] at (0,0) {B};

\end{axis}
\end{tikzpicture}
\end{center}

\item The basic direction associated with $x_3$ is
  \[
d_3 = \cvect{0 \\ 0 \\ d_B \\ 1},
\]
where
\[
d_B = -B^{-1}A_4 = -1,
\]
as the variable $x_3$ is associated with column 4 of $A$,
so that 
  \[
d_3 = \cvect{0 \\ 0 \\ -1 \\ 1}.
\]
Performing a step $\alpha$ in this direction leads to the point
\[
\cvect{0 \\ 0 \\ 1 \\ 0} + \alpha \cvect{0 \\ 0 \\ -1 \\ 1} =
\cvect{0 \\ 0 \\ 1-\alpha \\ \alpha}.
\]
The point is feasible for any value of $\alpha$ such that $0 \leq \alpha \leq 1$, so that the
direction is feasible.
In the original variables, with $x_1=x_1^+-x_1^-$, it corresponds to
a direction
\[
\cvect{0-0 \\ -1}=
\cvect{0 \\ -1},
\]
as illustrated on the following figure:
\begin{center}
\begin{tikzpicture}[scale=1.2]
  \begin{axis}[ xlabel={$x_1$}, ylabel={$x_2$}
      ,axis on top=true
      ,axis equal
      ,axis lines=middle
      ,samples=41
      ,xtick = {1}
      ,ytick = {1}
      ,domain=0:1.5
      ,xmin=-1.3,xmax=1.3
      ,ymin=-0.1,ymax=1.2
,legend pos=south east
]
\addplot [thick,color=gray!20,fill=gray!20, 
                    fill opacity=0.05]coordinates {
  ( -1.3,0) 
  ( -1.3,1)
  ( 1.3,1)
  (1.3,0)
};
\addplot[color=black, thin,domain=-1.2:1.2] {0}
node[pos=0.2,sloped,above] {\tiny $x_2=0$};
\addplot[color=black,thin,domain=-1.2:1.2] {1}
node[pos=0.2,sloped,below] {\tiny$x_2=1$};

\draw[-latex,very thick] (0,1) -- (0,0) node[pos=0.5,above,sloped] {$d_3$} ;
%\node [circle,fill=black,draw,scale=0.2,label=above right:$A$] at (0,1) {A};
%\node [circle,fill=black,draw,scale=0.2,label=below right:$B$] at (0,0) {B};

\end{axis}
\end{tikzpicture}
\end{center}

\end{enumerate}

